%!TeX root=../gardentop.tex
\chapter{»I Won't!« said Mary} 
	
\begin{figure}[t!]
\centering
\includegraphics[width=\linewidth]{june}
\end{figure}

	\lettrine[lines=6]{T}{hey} found a great deal to do that morning and Mary was late in returning to the house and was also in such a hurry to get back to her work that she quite forgot Colin until the last moment.

\zz
»Tell Colin that I can't come and see him yet,« she said to Martha. »I'm very busy in the garden.«

Martha looked rather frightened.

»Eh! Miss Mary,« she said, »it may put him all out of humor when I tell him that.«

But Mary was not as afraid of him as other people were and she was not a self-sacrificing person.

»I can't stay,« she answered. »Dickon's waiting for me;« and she ran away.

The afternoon was even lovelier and busier than the morning had been. Already nearly all the weeds were cleared out of the garden and most of the roses and trees had been pruned or dug about. Dickon had brought a spade of his own and he had taught Mary to use all her tools, so that by this time it was plain that though the lovely wild place was not likely to become a »gardener's garden« it would be a wilderness of growing things before the springtime was over.

»There'll be apple blossoms an' cherry blossoms overhead,« Dickon said, working away with all his might. »An' there'll be peach an' plum trees in bloom against th' walls, an' th' grass'll be a carpet o' flowers.«

The little fox and the rook were as happy and busy as they were, and the robin and his mate flew backward and forward like tiny streaks of lightning. Sometimes the rook flapped his black wings and soared away over the tree-tops in the park. Each time he came back and perched near Dickon and cawed several times as if he were relating his adventures, and Dickon talked to him just as he had talked to the robin. Once when Dickon was so busy that he did not answer him at first, Soot flew on to his shoulders and gently tweaked his ear with his large beak. When Mary wanted to rest a little Dickon sat down with her under a tree and once he took his pipe out of his pocket and played the soft strange little notes and two squirrels appeared on the wall and looked and listened.

»Tha's a good bit stronger than tha' was,« Dickon said, looking at her as she was digging. »Tha's beginning to look different, for sure.«

Mary was glowing with exercise and good spirits.

»I'm getting fatter and fatter every day,« she said quite exultantly. »Mrs~Medlock will have to get me some bigger dresses. Martha says my hair is growing thicker. It isn't so flat and stringy.«

The sun was beginning to set and sending deep gold-coloured rays slanting under the trees when they parted.

»It'll be fine to-morrow,« said Dickon. »I'll be at work by sunrise.«

»So will I,« said Mary.

She ran back to the house as quickly as her feet would carry her. She wanted to tell Colin about Dickon's fox cub and the rook and about what the springtime had been doing. She felt sure he would like to hear. So it was not very pleasant when she opened the door of her room, to see Martha standing waiting for her with a doleful face.

»What is the matter?« she asked. »What did Colin say when you told him I couldn't come?«

»Eh!« said Martha, »I wish tha'd gone. He was nigh goin' into one o' his tantrums. There's been a nice to do all afternoon to keep him quiet. He would watch the clock all th' time.«

Mary's lips pinched themselves together. She was no more used to considering other people than Colin was and she saw no reason why an ill-tempered boy should interfere with the thing she liked best. She knew nothing about the pitifulness of people who had been ill and nervous and who did not know that they could control their tempers and need not make other people ill and nervous, too. When she had had a headache in India she had done her best to see that everybody else also had a headache or something quite as bad. And she felt she was quite right; but of course now she felt that Colin was quite wrong.

He was not on his sofa when she went into his room. He was lying flat on his back in bed and he did not turn his head toward her as she came in. This was a bad beginning and Mary marched up to him with her stiff manner.

»Why didn't you get up?« she said.

»I did get up this morning when I thought you were coming,« he answered, without looking at her. »I made them put me back in bed this afternoon. My back ached and my head ached and I was tired. Why didn't you come?«

»I was working in the garden with Dickon,« said Mary.

Colin frowned and condescended to look at her.

»I won't let that boy come here if you go and stay with him instead of coming to talk to me,« he said.

Mary flew into a fine passion. She could fly into a passion without making a noise. She just grew sour and obstinate and did not care what happened.

»If you send Dickon away, I'll never come into this room again!« she retorted.

»You'll have to if I want you,« said Colin.

»I won't!« said Mary.

»I'll make you,« said Colin, »They shall drag you in.«

»Shall they, Mr~Rajah!« said Mary fiercely. »They may drag me in but they can't make me talk when they get me here. I'll sit and clench my teeth and never tell you one thing. I won't even look at you. I'll stare at the floor!«

They were a nice agreeable pair as they glared at each other. If they had been two little street boys they would have sprung at each other and had a rough-and-tumble fight. As it was, they did the next thing to it.

»You are a selfish thing!« cried Colin.

»What are you?« said Mary. »Selfish people always say that. Any one is selfish who doesn't do what they want. You're more selfish than I am. You're the most selfish boy I ever saw.«

»I'm not!« snapped Colin. »I'm not as selfish as your fine Dickon is! He keeps you playing in the dirt when he knows I am all by myself. He's selfish, if you like!«

Mary's eyes flashed fire.

»He's nicer than any other boy that ever lived!« she said. »He's—he's like an angel!« It might sound rather silly to say that but she did not care.

»A nice angel!« Colin sneered ferociously. »He's a common cottage boy off the moor!«

»He's better than a common Rajah!« retorted Mary. »He's a thousand times better!«

Because she was the stronger of the two she was beginning to get the better of him. The truth was that he had never had a fight with any one like himself in his life and, upon the whole, it was rather good for him, though neither he nor Mary knew anything about that. He turned his head on his pillow and shut his eyes and a big tear was squeezed out and ran down his cheek. He was beginning to feel pathetic and sorry for himself—not for any one else.

»I'm not as selfish as you, because I'm always ill, and I'm sure there is a lump coming on my back,« he said. »And I am going to die besides.«

»You're not!« contradicted Mary unsympathetically.

He opened his eyes quite wide with indignation. He had never heard such a thing said before. He was at once furious and slightly pleased, if a person could be both at the same time.

»I'm not?« he cried. »I am! You know I am! Everybody says so.«

»I don't believe it!« said Mary sourly. »You just say that to make people sorry. I believe you're proud of it. I don't believe it! If you were a nice boy it might be true—but you're too nasty!«

In spite of his invalid back Colin sat up in bed in quite a healthy rage.

»Get out of the room!« he shouted and he caught hold of his pillow and threw it at her. He was not strong enough to throw it far and it only fell at her feet, but Mary's face looked as pinched as a nutcracker.

»I'm going,« she said. »And I won't come back!«

She walked to the door and when she reached it she turned round and spoke again.

»I was going to tell you all sorts of nice things,« she said. »Dickon brought his fox and his rook and I was going to tell you all about them. Now I won't tell you a single thing!«

She marched out of the door and closed it behind her, and there to her great astonishment she found the trained nurse standing as if she had been listening and, more amazing still—she was laughing. She was a big handsome young woman who ought not to have been a trained nurse at all, as she could not bear invalids and she was always making excuses to leave Colin to Martha or any one else who would take her place. Mary had never liked her, and she simply stood and gazed up at her as she stood giggling into her handkerchief.

»What are you laughing at?« she asked her.

»At you two young ones,« said the nurse. »It's the best thing that could happen to the sickly pampered thing to have some one to stand up to him that's as spoiled as himself;« and she laughed into her handkerchief again. »If he'd had a young vixen of a sister to fight with it would have been the saving of him.«

»Is he going to die?«

»I don't know and I don't care,« said the nurse. »Hysterics and temper are half what ails him.«

»What are hysterics?« asked Mary.

»You'll find out if you work him into a tantrum after this—but at any rate you've given him something to have hysterics about, and I'm glad of it.«

Mary went back to her room not feeling at all as she had felt when she had come in from the garden. She was cross and disappointed but not at all sorry for Colin. She had looked forward to telling him a great many things and she had meant to try to make up her mind whether it would be safe to trust him with the great secret. She had been beginning to think it would be, but now she had changed her mind entirely. She would never tell him and he could stay in his room and never get any fresh air and die if he liked! It would serve him right! She felt so sour and unrelenting that for a few minutes she almost forgot about Dickon and the green veil creeping over the world and the soft wind blowing down from the moor.

Martha was waiting for her and the trouble in her face had been temporarily replaced by interest and curiosity. There was a wooden box on the table and its cover had been removed and revealed that it was full of neat packages.

»Mr~Craven sent it to you,« said Martha. »It looks as if it had picture-books in it.«

Mary remembered what he had asked her the day she had gone to his room. »Do you want anything—dolls—toys—books?« She opened the package wondering if he had sent a doll, and also wondering what she should do with it if he had. But he had not sent one. There were several beautiful books such as Colin had, and two of them were about gardens and were full of pictures. There were two or three games and there was a beautiful little writing-case with a gold monogram on it and a gold pen and inkstand.

Everything was so nice that her pleasure began to crowd her anger out of her mind. She had not expected him to remember her at all and her hard little heart grew quite warm.

»I can write better than I can print,« she said, »and the first thing I shall write with that pen will be a letter to tell him I am much obliged.«

If she had been friends with Colin she would have run to show him her presents at once, and they would have looked at the pictures and read some of the gardening books and perhaps tried playing the games, and he would have enjoyed himself so much he would never once have thought he was going to die or have put his hand on his spine to see if there was a lump coming. He had a way of doing that which she could not bear. It gave her an uncomfortable frightened feeling because he always looked so frightened himself. He said that if he felt even quite a little lump some day he should know his hunch had begun to grow. Something he had heard Mrs~Medlock whispering to the nurse had given him the idea and he had thought over it in secret until it was quite firmly fixed in his mind. Mrs~Medlock had said his father's back had begun to show its crookedness in that way when he was a child. He had never told any one but Mary that most of his »tantrums« as they called them grew out of his hysterical hidden fear. Mary had been sorry for him when he had told her.

»He always began to think about it when he was cross or tired,« she said to herself. »And he has been cross to-day. Perhaps—perhaps he has been thinking about it all afternoon.«

She stood still, looking down at the carpet and thinking.

»I said I would never go back again—« she hesitated, knitting her brows—»but perhaps, just perhaps, I will go and see—if he wants me—in the morning. Perhaps he'll try to throw his pillow at me again, but—I think—I'll go.«

